\chapter{Variables and Scope}

\section{Declaration with \texttt{let}}

Every variable is declared with \lstinline{let}. The declaration binds a
name to a value:

\begin{lstlisting}
let x    = 10
let name = "Hayashi"
let active = true
\end{lstlisting}

\lstinline{let} without an initial value does not exist in \hayashi{} —
every variable is initialised at the point of declaration.

\subsection*{Redeclaration}

A variable may be redeclared with \lstinline{let} in the same scope. The
new value replaces the previous one without affecting other variables that
point to the old value:

\begin{lstlisting}
let x = 5
let y = x      // y points to 5
let x = 99     // x is now 99; y remains 5
display y      // 5
\end{lstlisting}

\section{Constants with \texttt{const}}

\lstinline{const} declares a name that cannot be redeclared or reassigned
in the same scope. Useful for configuration values or fixed parameters:

\begin{lstlisting}
const ALPHA   = 0.05
const PROJECT = "my_lab"
\end{lstlisting}

Attempting to assign to a constant produces an error:

\begin{lstlisting}
const N = 100
N = 200    // error: cannot assign to constant 'N'
\end{lstlisting}

\begin{nota}
  \lstinline{const} in \hayashi{} is evaluated at runtime, not at compile
  time. The distinction from \lstinline{let} is purely one of binding
  mutability, not performance.
\end{nota}

\section{Assignment}

After declaration with \lstinline{let}, a variable can be reassigned with
\lstinline{=}:

\begin{lstlisting}
let counter = 0
counter = counter + 1
counter = counter + 1
display counter    // 2
\end{lstlisting}

Compound assignment operators:

\begin{lstlisting}
let n = 10
n += 5     // n = 15
n -= 3     // n = 12
n *= 2     // n = 24
n /= 4     // n = 6
\end{lstlisting}

\section{Scopes}

Scopes are delimited by \lstinline|{ }| blocks. Variables declared inside
a block exist only until the end of that block:

\begin{lstlisting}[caption={Variables local to a block}]
let x = 1

{
  let y = 2
  display x    // 1  — x is visible here
  display y    // 2
}

display x    // 1
display y    // error: undefined variable 'y'
\end{lstlisting}

\subsection*{Scopes in functions}

Each function call creates a new scope. Parameters and local variables are
destroyed when the function returns:

\begin{lstlisting}
fn sum(a, b) {
  let result = a + b
  return result
}

display sum(3, 4)    // 7
display result       // error: undefined variable 'result'
\end{lstlisting}

\subsection*{Outer scope capture (closures)}

Anonymous functions (closures) capture variables from the scope where they
were defined:

\begin{lstlisting}
let base = 100

let add = fn(x) {
  return x + base    // captures 'base' from outer scope
}

display add(20)    // 120
\end{lstlisting}

\section{Special variables}

\subsection*{\texttt{\_n} and \texttt{\_N} in data contexts}

Inside \lstinline{generate} and \lstinline{mutate}, two implicit variables
are available:

\begin{itemize}
  \item \lstinline{_n}: index of the current row (starting at 1)
  \item \lstinline{_N}: total number of rows in the DataFrame
\end{itemize}

\begin{lstlisting}[caption={Using \_n and \_N}]
input x
  10  20  30
end

generate df id     = _n       // [1, 2, 3]
generate df weight = 1/_N     // [0.333, 0.333, 0.333]
\end{lstlisting}

\subsection*{\texttt{\_} (discard)}

The identifier \lstinline{_} discards a value without creating a binding:

\begin{lstlisting}
for _ in 1..5 {
  display "iteration"
}
\end{lstlisting}
